\section{Fano多様体のtensor根と曲面のPoisson方向}
\label{sec:surface}
$X$を滑らかな複素Fano多様体とする．HKR分解は
\begin{equation}\label{eq:hkr}
 \HH^2(X)=U\oplus P,\qquad U=H^1(X,T_X),\quad P=H^0(X,\wedge^2T_X)
\end{equation}
となる．$H^2(\OO_X)=0$を使った．曲面では$P=H^0(-K_X)$で，全双ベクトルはPoissonである．
高次元の双ベクトルは一次変形を定めるが，二次以上にはJacobi障害があり得る．$P$の作用は$\wedge^2T_X$の\emph{自然な線型化}で取り，
反標準写像の射影作用だけで置き換えない．

\begin{lemma}[三角作用]\label{lem:triangular}
$s=L\otimes f^*$の$\HH^2$作用は，\eqref{eq:hkr}について
\[
 \alpha(u,\pi)=(f^*u+B_L\pi,\ f^*\pi)
\]
という形である．ここで$B_L:P\to U$はある線型写像である．
\end{lemma}
\begin{proof}
線束のtensorは局所的に恒等Morita同値なので，局所的な非可換変形の記号である
大域Poisson成分を変えない．一方，$H^2(\OO_X)=0$により$L$は任意の一次可換変形へ持ち上がり，
そのtensor自己同値は同じ可換変形上にある．従ってtensorは$U$を恒等に保つ．
$f^*$は通常のHKR分解と互換である．以上から表示を得る．
\end{proof}

\begin{theorem}[通常の指数のtensor根は一次変形を制限しない]\label{thm:tensor-root}
滑らかな複素Fano多様体で$L^{\otimes r}\simeq\omega_X^{-1}$なら，
$T_L=(-)\otimes L$は$\HH^2(X)$に恒等に作用する．従って任意の一次Abel変形に
$T_L$はexact自己同値として持ち上がる．
さらに$s=L\otimes f^*$が根であるとき，その一次持上げ条件は単に
$f_{\HH}(u)=u$である．相対Serre仮定の下では，これらは根の冪関係も満たす持上げになる．
\end{theorem}
\begin{proof}
補題\ref{lem:triangular}で$f=1$とすれば，$N=(T_L)_{\HH}-1$は$P$から$U$へ写し，$U$を消すので$N^2=0$．
$(T_L)^r=\tau$と命題\ref{prop:finite-hh}より$(1+N)^r=1$だから$rN=0$，従って$N=0$である．
定理\ref{thm:fixed}と系\ref{cor:relative-root}を適用する．最後の条件では
$s_{\HH}=(T_L)_{\HH}f_{\HH}=f_{\HH}$である．
\end{proof}

\begin{proposition}[不変Poisson方向の補償]\label{prop:invariant-poisson}
$s^r=\tau$なら，以下は完全列である：
\begin{equation}\label{eq:invariants}
 0\longrightarrow U^{f^*}\longrightarrow\HH^2(X)^{s_{\HH}}
 \longrightarrow P^{f^*}\longrightarrow0.
\end{equation}
さらに$f^*|_U=1$なら，HKRの\emph{元の分解}において
\begin{equation}\label{eq:trivial-U}
 \HH^2(X)^{s_{\HH}}=U\oplus P^{f^*}.
\end{equation}
\end{proposition}
\begin{proof}
$0\to U\to\HH^2\to P\to0$は命題\ref{prop:finite-hh}により有限巡回群の表現の完全列である．
標数零の不変部分関手は平均化によりexactなので\eqref{eq:invariants}を得る．
$f^*|_U=1$とし$\pi\in P^{f^*}$を取ると，補題\ref{lem:triangular}より
$\alpha^r(0,\pi)=(rB_L\pi,\pi)$．$\alpha^r=1$だから$B_L\pi=0$である．
従って$(0,\pi)$自身が不変となり，\eqref{eq:trivial-U}が従う．
\end{proof}

\begin{remark}[Picard re-markingとの接続]\label{rem:picard}
局所積を$ab+\varepsilon\pi(da,db)/2$と規格化し，線束の遷移関数を$g_{ij}$とすると，
そのstar共役は$\id+\varepsilon\pi^\sharp(d\log g_{ij})$である．
従ってtensorによるPoisson成分から可換成分への混合は，向きの規約を固定すれば
\[
 \delta_\pi(L)=H^1(\pi^\sharp)c_1(L)\in H^1(T_X)
\]
で与えられる．符号を反転してもその核は変わらない．
これは別研究\cite{Marking}のmarking忘却の障害と同じ一次機構であり，
\emph{線束を対象として持ち上げること}と\emph{tensorを固定変形の自己同値として持ち上げること}は異なる．
Bursztyn--Dolgushev--Waldmann\cite[Theorem 3.11]{BDW}は滑らかな実多様体上のstar productで
Picard--Morita作用を$B$-field作用と同定する．その$C^\infty$形式理論を，本稿の射影代数的変形へ無条件に転用しない．
\end{remark}

\begin{proposition}[圏自身の形式的非障害性]\label{prop:hh3}
滑らかな複素del Pezzo曲面では$\HH^3(X)=0$である．
従ってLowen--Van den BerghのAbel変形関手では，小拡大に関する障害が消える．
これ単独では根の持上げを意味しないが，$\HH^1=0$のとき定理\ref{thm:formal-roots}と組み合わせられる．
相対Fano性は命題\ref{prop:relative-ample}の別の議論を使い，代数化はなお含まない．
\end{proposition}
\begin{proof}
HKRにより$\HH^3=H^1(-K_X)\oplus H^2(T_X)$．
$-K_X=K_X+(-2K_X)$とKodaira消滅から第一項は零．
Serre双対性から第二項の双対は$H^0(\Omega_X^1\otimes K_X)$である．
非零反標準切断を掛ける単射$\Omega_X^1\otimes K_X\hookrightarrow\Omega_X^1$と
合理曲面の$H^0(\Omega_X^1)=0$によりこれも零．
変形の障害は\cite[Theorem 3.1]{LVdB}を用いる．
\end{proof}

